% ── Abstract ─────────────────────────────────────────────────────────
% Length  : 150–250 words
% Structure (fill AFTER experiments):
%   [P] Problem    — fixed thresholds fail in mixed traffic
%   [G] Gap        — no system adapts perception for mixed + dual-task
%   [M] Method     — C1 context estimator, C2 CTRV model, C3 reliability score
%   [R] Results    — numbers: MOTA↑, FP rate↓, accident F1↑
% ─────────────────────────────────────────────────────────────────────
% TODO: write after Section V (Results) is complete.

Urban traffic surveillance in motorcycle-dominant environments poses
unique challenges for perception systems.
%
[P] Fixed detection and tracking thresholds lead to high false-positive
accident alerts and degraded multi-object tracking during dense weaving
maneuvers, directly reducing the reliability of downstream safety and
control applications.
%
[G] Existing systems address either highway accident detection or
multi-agent traffic control, but none adapt perception parameters to
real-time traffic context, particularly for mixed-traffic scenes where
motorcycles constitute over 90\% of vehicles.
%
[M] We propose a context-aware adaptive perception pipeline comprising
three contributions: (C1)~a lightweight context estimator that adjusts
YOLO confidence thresholds and ByteTrack buffer sizes based on traffic
density, time-of-day, and lighting; (C2)~a class-conditioned motion
model replacing linear Kalman prediction with the constant turn-rate and
velocity (CTRV) model for motorcycles; and (C3)~a dual-task track
reliability score that simultaneously optimizes precision for accident
alerting and recall for MAPPO-based signal control observations.
%
[R] Evaluated on a self-collected Vietnamese urban dataset
(TODO: numbers here --- e.g., our method achieves MOTA of XX\%,
IDF1 of XX\%, and reduces false-positive alert rate by XX\%
compared to the vanilla ByteTrack baseline.)
